% Paper skeleton: static count-one ownership and in-place array updates in Lean 4.
%
% Build note: the Lean blocks use minted, so compile with
%   lualatex -shell-escape ownership-skeleton.tex
% and a monospace font that covers Lean's Unicode glyphs (adjust \setmonofont).
% The \prose{...} markers are visible placeholders (todonotes); remove the
% package/macro when the prose is written. No prose is supplied.
%
% Theorem statements are the repository's actual signatures, with premises
% abbreviated where marked by "-- …". Bibliographic details to be verified.
%
% Build note on this host: the installed Debian TeX Live lacks texlive-luatex
% (luaotfload), so fontspec cannot load fonts. Without root, fetch and extract
% that one package:
%   apt-get download texlive-luatex
%   dpkg-deb -x texlive-luatex_*.deb /tmp/tl-luatex/root
%   export TEXMFLOCAL=/tmp/tl-luatex/root/usr/share/texlive/texmf-dist
%   mktexlsr "$TEXMFLOCAL"
% then
%   lualatex -shell-escape ownership-skeleton.tex
% (minted needs pygments). If a Unicode mono font is unavailable, replace the
% fontspec/minted block with the listings fallback commented below; the Lean
% glyphs then need literate mappings or a Unicode-capable mono font.

\documentclass[11pt]{article}

\usepackage{fontspec}
\setmonofont{DejaVu Sans Mono}[Scale=0.85]
\usepackage{minted}
\setminted{breaklines=true, fontsize=\small, frame=lines, framesep=2mm}
% Fallback without fontspec/minted (also needs literate mappings for ∀ → ∧ ∃):
% \usepackage[T1]{fontenc}\usepackage{lmodern}\usepackage{listings}
% \lstset{basicstyle=\ttfamily\small, breaklines=true}
\usepackage{booktabs}
\usepackage{hyperref}
\usepackage{cleveref}
\usepackage{todonotes}
\newcommand{\prose}[1]{\todo[inline, color=yellow!20, size=\small]{#1}}

\newtheorem{claim}{Claim}
\newtheorem{definition}{Definition}
\newtheorem{nonclaim}{Non-claim}

\title{Static Count-One Ownership for In-Place Array Updates in Lean~4}
\author{[Author]}
\date{[Date]}

\begin{document}
\maketitle

\begin{abstract}
\prose{One paragraph: the problem (Lean arrays copy when shared), the claim
(an ownership discipline with a count-one invariant yields in-place updates for
the recognised emitted write), and the boundary (what is proved, what is
measured, what is trusted).}
\end{abstract}

\section*{Argument map}
% ARG: the paper's spine; each row is one step of the argument and the theorem
% that carries it. Prose should introduce each row and then hand over to the
% claims.
\begin{center}
\begin{tabular}{@{}lll@{}}
\toprule
Step & Claim & Carrier \\
\midrule
1 & the checker is sound: accepted writes are exclusive & \Cref{thm:check-sound} \\
2 & the entry invariant supplies count one and separation & \Cref{thm:entry-inv} \\
3 & the emitted write prefix is recognised and corresponds & \Cref{thm:prefix-sound} \\
4 & the recognised write lands in place & \Cref{thm:prefix-in-place} \\
5 & the nested path corresponds without a callee hypothesis & \Cref{thm:path-sound} \\
6 & the extraction pair returns to count one & \Cref{thm:pair} \\
7 & the put-back's exclusive branch is derived & \Cref{thm:putback} \\
8 & the runtime model cannot abort under the correspondence & \Cref{thm:no-abort} \\
9 & the observable behaviour agrees on a corpus & \Cref{sec:empirical} \\
\bottomrule
\end{tabular}
\end{center}

\section{Introduction}
\prose{Motivate the gap: a purely functional language's updates copy when the
array is shared; the paper asks whether a static discipline can guarantee the
in-place branch. State the two claims (model-level theorem; runtime-model
simulation) and the empirical boundary. Cite \cite{lean4,perceus,countingbeans}.}
\prose{List the contributions, each pointing at a claim below, and state the
non-claims early (no C-semantics proof, no M7 source correspondence).}

\section{Setting: Lean 4, LCNF, and reference counting}
\prose{Describe the pinned toolchain, LCNF, the runtime's \texttt{lean\_is\_exclusive}
and \texttt{lean\_array\_uset}, and the representation of arrays. Cite
\cite{lean4,leanruntime,perceus}.}
\prose{State the problem precisely: the runtime copies unless the count is one and
the object is non-persistent; a source-level discipline must arrange that.}

\section{The ownership discipline}
\prose{Introduce the model heap, values, counts, exclusivity, and the checker as
a proof obligation, not a runtime test. The key point: count one is a ghost
predicate enforced by the checker and preserved by the invariant.}

\begin{definition}[Invariant]\label{def:inv}
\begin{minted}{lean}
structure Inv (H : Heap) (env : Env) (c : Ctx) : Prop where
  rc_eq       : ∀ r, H.rc r = c.mult r
  owned       : ∀ x r, c.us x = .owned → c.refOf x = some r →
                  env x = some (.ref r) ∧ 1 ≤ c.mult r
  borrow      : ∀ x r, c.us x = .borrow → c.refOf x = some r →
                  env x = some (.ref r) ∧ H.cells r ≠ none
  scalar      : ∀ x, c.us x = .scalar → ∃ n, env x = some (.scalar n)
  nonPersistent : ∀ r o, H.cells r = some o → o.persistent = false

def Exclusive (H : Heap) (r : Ref) : Prop :=
  ∃ o, H.cells r = some o ∧ o.rc = 1 ∧ o.persistent = false
\end{minted}
\end{definition}

\begin{claim}[Checker soundness]\label{thm:check-sound}
\begin{minted}{lean}
theorem check_sound {c : Code} {ctx ctx' : Ctx} {H : Heap} {env : Env}
    (hc : check ctx c = some ctx') (hi : Inv H env ctx) :
    ∃ H' env' v ws, run c env H = .done H' env' v ws ∧
      (∀ w ∈ ws, w.2 = true) ∧ Inv H' env' ctx'
\end{minted}
\end{claim}

\begin{claim}[Entry invariant]\label{thm:entry-inv}
\begin{minted}{lean}
structure EntryInv (H : Heap) (env : Env) (c : Ctx) (root : Ref)
    (client : List Ref) : Prop extends Inv H env c where
  uniqueTree  : UniqueTree H root                 -- rc = 1 on the footprint
  separated   : ClientSeparated H root client
  complete    : RootsComplete c H root

theorem checkN_sound_of_entryInv -- … premises abbreviated
    (hEntry : EntryInv H env c root client) (hag : Agrees H env code)
    (hc : checkN c code = some c') :
    ∃ H' env' v ws, runN code env H = .done H' env' v ws ∧
      (∀ w ∈ ws, w.2 = true) ∧ Inv H' env' c'
\end{minted}
\end{claim}

\prose{Path-level agreement and annotation reachability replace per-client
premises: \texttt{agrees\_of\_path}, \texttt{annotationsIn\_of\_path}, the
transitivity lemmas, and \texttt{path\_run\_safe}/\texttt{client\_retention}.
State what remains supplied.}

\section{Recognising the emitted write}
\prose{Explain the strict-recogniser approach: decode the actual LCNF value,
pin operation order, arities, operand identities and branch constructors, and
reject token-preserving near misses. The recogniser is structural; the contracts
A1--A12 (execP) and N1--N12 (nested) are named assumptions.}
\prose{Describe the three recognised shapes: \texttt{execP}'s \texttt{PTree.set},
\texttt{setAt}'s \texttt{Path.index}, and \texttt{setAt}'s
\texttt{Path.left}/\texttt{Path.right} pair arms (with joins). Cite the capture
artifact and the near-miss counts (11, 10, 28).}

\section{Correspondence to the ownership model}
\prose{This is the semantic bridge: the recognised fragment's restricted
semantics equals the model run. Emphasise that the recognition hypothesis is
load-bearing (\texttt{execAST} returns nothing unless the code decodes).}

\begin{claim}[Prefix correspondence]\label{thm:prefix-sound}
\begin{minted}{lean}
theorem project_sound (assign : FVarId → Var) (p : Prefix)
    (code : Code .impure) (hproj : projectPrefix code = some p) -- … premises
    (hag : Agrees assign p s env H root i v) (hcell : H.cells root = some o)
    (hin : i < o.fields.size) :
    PrefixOutcome.ofExec assign p (execAST assign p code s) =
      some (PrefixOutcome.ofModel (Ownership.run (p.toModel assign) env H))
\end{minted}
\end{claim}

\begin{claim}[The write is in place]\label{thm:prefix-in-place}
\begin{minted}{lean}
theorem recognized_write_in_place -- … premises abbreviated
    (hEntry : EntryInv H env (Runner.initCtx root) root client)
    (hproj : projectPrefix code = some p) (hl : Layout assign p) -- …
    (hcell : H.cells root = some o) (hin : i < o.fields.size) :
    PrefixOutcome.ofExec assign p (execAST assign p code s) =
        some (inPlaceOutcome root i v o H) ∧
      PrefixOutcome.ofModel (Ownership.run (p.toModel assign) env H) =
        inPlaceOutcome root i v o H
\end{minted}
\end{claim}

\prose{The nested path: the model builds take-out/call/put-back per level; the
finite-path induction discharges the callee contract from the tail, leaving the
leaf contract and N10. The model bound: the flat environment cannot express the
emitted per-invocation \texttt{oproj[1]} index; record it as a limitation.}

\begin{claim}[Nested finite path]\label{thm:path-sound}
\begin{minted}{lean}
theorem index_path_simulates_model (declName takeOut : Name) -- … premises
    (hproj : projectNested declName takeOut code = some p) -- …
    (build : …) (ls : List Level) :
    NestedOutcome.ofExec assign p
        (execNestedAST declName takeOut assign body p code s) =
      NestedOutcome.ofModel (Nested.runN (build path value body ls) mEnv H) ∧
    ∃ H' env' ws, Nested.runN (build path value body ls) mEnv H =
        .done H' env' (.ref rA) ws ∧ ∀ w ∈ ws, w.2 = true
\end{minted}
\end{claim}

\begin{minted}{lean}
theorem calls_build (path value : Var) (body : NCode) (ls : List Level) :
    calls (build path value body ls) = calls body + ls.length
theorem callBody_build -- (l :: rest), rfl:
    callBody? (build path value body (l :: rest)) = some (build path value body rest)
\end{minted}

\section{Counted extraction, call, and restoration}
\prose{The pair's temporary handles: the take-out transfers the parent and child
into the fresh pair, the projections increment, and freeing the unique pair
releases both edges, so the boundary counts equal the entry counts. The negative
controls are what make the claim non-vacuous.}

\begin{claim}[Pair accounting]\label{thm:pair}
\begin{minted}{lean}
theorem pair_accounting {H : Heap} {p c : Ref} {i : Nat}
    (hw : WellFormedEntry H p c i) :
    let B := protocol H p c i
    B.2.rc p = 1 ∧ B.2.rc c = 1 ∧ B.2.cells p ≠ none ∧ B.2.cells c ≠ none ∧
      B.2.cells B.1 = none ∧ Heap.Exclusive B.2 p ∧ Heap.Exclusive B.2 c ∧
      B.2.rc p = H.rc p ∧ B.2.rc c = H.rc c

theorem no_free_components_not_exclusive -- negative control
    (hw : WellFormedEntry H p c i) :
    let B := protocolNoFree H p c i
    B.2.rc p = 2 ∧ B.2.rc c = 2 ∧ ¬ Heap.Exclusive B.2 p
\end{minted}
\end{claim}

\prose{Put-back: two dispositions, the exclusive branch derived from the
counted-pair boundary and the child-call contract. Functional restoration and
count preservation are separate. The child-call boundary is derived (confinement
from \texttt{callAt\_frame\_confined}), not assumed; the returned-handle counts
remain premises.}

\begin{claim}[Put-back exclusivity]\label{thm:putback}
\begin{minted}{lean}
theorem putBack_exclusive_of_wellFormed -- … premises abbreviated
    (hw : CountedPair.WellFormedEntry H p c i)
    (hc : ChildCall Hin Hout env ctx state path x body p c) :
    putBack H p i c = some (p, mutate H p o i c) ∧ -- … poststate
      Heap.Exclusive (mutate H p o i c) p ∧ H.rc p = 1
\end{minted}
\end{claim}

\section{Runtime model and abort impossibility}
\prose{Define the independent runtime state, its exclusivity predicate and the
two array-set dispositions; the structural correspondence preserves refs, cells,
counts, persistence, fields and the allocator cursor. The abort branch is
unreachable under the correspondence and \texttt{EntryInv}; a guarded allocator
refuses a dirty cursor. Name the fidelity boundary: the model is the documented
meaning of the primitives, not a proof of the C semantics.}

\begin{claim}[No abort]\label{thm:no-abort}
\begin{minted}{lean}
structure RHeap where
  cells : Ref → Option RObj
  next  : Ref

def RHeap.isExclusive (R : RHeap) (r : Ref) : Bool := -- rc = 1 and not persistent

def checkExclusive (R : RHeap) (r : Ref) : CheckOutcome := -- .ok r or .abort r

def arraySet (R : RHeap) (r : Ref) (i : Nat) (v : Val) : Option SetOutcome := -- inPlace or copied

structure Corr (H : Heap) (R : RHeap) : Prop where -- refs, cells, rc, persistent, fields, next

theorem checkExclusive_no_abort {H : Heap} {R : RHeap} {r : Ref}
    (hC : Corr H R) (hx : H.Exclusive r) :
    checkExclusive R r = .ok r ∧ ∀ r', checkExclusive R r ≠ .abort r'

theorem setAt_inPlace_of_entryInv -- … premises abbreviated
    (hEntry : EntryInv H env (Runner.initCtx root) root client)
    (hC : Corr H R) (hcell : H.cells root = some o)
    (hR : R.cells root = some o') (hin : i < o.fields.size) :
    ∃ H', setAt R root i (.scalar vn) = some (.inPlace H' root) ∧
      Corr (inPlaceOutcome root i vn o H).heap H'
\end{minted}
\end{claim}

\begin{claim}[Emitted-semantics abort impossibility]\label{thm:prefix-no-abort}
\begin{minted}{lean}
def PrefixAborts (out : Option PrefixOutcome) : Prop :=
  ∃ o, out = some o ∧ o.fault ≠ none

theorem prefix_abort_impossible -- … the recognised-prefix premises
    (hproj : projectPrefix code = some p) -- …
    (hEntry : EntryInv H env (Runner.initCtx root) root client) :
    ¬ PrefixAborts (PrefixOutcome.ofExec assign p (execAST assign p code s))
\end{minted}
\end{claim}

\section{Empirical alignment}\label{sec:empirical}
\prose{The compiled corpus and the gate: the model's predicates against
\texttt{lean\_is\_exclusive} and address identity, with the check armed and
unarmed. Report the context finding (closed pure arrays are lifted to
\texttt{\_\_\_closed}/lean\_obj\_once shared constants, so they copy) and the
differential corpus (unique, alias, drop-back-to-one, persistent, closure). State
that this is validation, not an exhaustive equivalence. Cite the pinned
artifact, IR-shape gate, generated C and runtime header.}

\begin{center}
\begin{tabular}{@{}lll@{}}
\toprule
Case & Predicted & Observed \\
\midrule
unique array & exclusive, in place & agree \\
live alias (count two) & shared, copy & agree \\
alias released before the write & exclusive, in place & agree \\
persistent literal & shared, copy & agree \\
closure-held array & shared, copy & agree \\
\bottomrule
\end{tabular}
\end{center}

\section{Limits and non-claims}
\begin{nonclaim}
\prose{Enumerate: no C-semantics refinement; A10 replay/backend identity; the
emitted recursive callee and continuation (N10); the model's flat-environment
index bound; address identity; free/cascade; concurrency and callbacks; M7
source correspondence; the leaf body contract and the remaining heap-shape
premises.}
\end{nonclaim}

\section{Related work}
\prose{Four clusters, one sentence each, citing the references:
reference counting and reuse \cite{perceus,countingbeans,koka}; linear and
uniqueness disciplines \cite{wadler-linear,linearhaskell,uniqueness,aliastypes,
ownershiptypes,cyclone}; program logics and ownership \cite{separationlogic,
iris,rustbelt,verus}; verified compilation and kernels \cite{compcert,cakeml,
lean4,lean4lean}.}

\section{Conclusion}
\prose{Restate what is proved (model-level and runtime-model) and what is
trusted (the primitives' C semantics and A10). Sketch the next step: a verified
runtime model or a backend gate with a stronger fidelity argument.}

\appendix
\section{Theorem map and artifacts}
\prose{Table of Lean name, file, role, and the claimed status (proved,
conditional, measured, trusted). Artifact pins, gate commands and the sibling
revision.}
\begin{center}
\begin{tabular}{@{}lll@{}}
\toprule
Lean name & File & Status \\
\midrule
\texttt{check\_sound} & \texttt{Ownership/Soundness.lean} & proved \\
\texttt{checkN\_sound\_of\_entryInv} & \texttt{Ownership/Entry.lean} & proved \\
\texttt{project\_sound} & \texttt{OwnershipLcnf/Semantics.lean} & proved (contracts) \\
\texttt{recognized\_write\_in\_place} & \texttt{OwnershipLcnf/Semantics.lean} & proved (contracts) \\
\texttt{index\_path\_simulates\_model} & \texttt{OwnershipLcnf/PathRecursion.lean} & proved (leaf/N10) \\
\texttt{pair\_accounting} & \texttt{Ownership/CountedPair.lean} & proved (restricted free) \\
\texttt{putBack\_exclusive\_of\_wellFormed} & \texttt{Ownership/PutBack.lean} & proved (returned-handle premises) \\
\texttt{checkExclusive\_no\_abort} & \texttt{OwnershipLcnf/RuntimeModel.lean} & proved (model) \\
\texttt{prefix\_abort\_impossible} & \texttt{OwnershipLcnf/NoAbort.lean} & proved (A6/A8) \\
\texttt{differential corpus} & \texttt{runtime-bridge} & measured \\
\bottomrule
\end{tabular}
\end{center}

\bibliographystyle{plain}
\bibliography{references}

\end{document}
